% ЗАКЛЮЧЕНИЕ, без номер.
\section*{ЗАКЛЮЧЕНИЕ}
\label{sec:conclusion}
\addcontentsline{toc}{section}{ЗАКЛЮЧЕНИЕ}

Проектът реализира клиентска библиотека на C++20, която говори пряко протоколите на PostgreSQL
и MySQL, с асинхронен пул, кеш на подготвените заявления и декодиране без копие. Изборът е
обоснован в Раздел~\ref{sec:analysis}: прякото говорене прави цената на достъпа наблюдаема.

\textbf{Предимства.} Липсата на междинни слоеве премахва копия, които иначе са скрити, и
измерването го потвърждава: декодирането от двоичен формат е 7,4 пъти по-евтино от текстовото,
а изборът на формат е достъпен само защото клиентът сам изпраща \code{Bind}. Кешът спестява 42
на сто от протоколните съобщения за едно действие. Асинхронният модел не изисква нишка на
връзка: всичко се изпълнява в един цикъл на събитията, без ключове за заключване, а проектът
се сглобява само с компилатор и CMake.

\textbf{Недостатъци и ограничения.} Проектът поема две протоколни реализации и риска от
разминаване при нови версии. Обхватът на протоколите не е пълен, съзнателно. Извън него
остават: за PostgreSQL режимът \code{COPY},
\code{LISTEN}/\code{NOTIFY}, отмяната на изпълняващо се заявление по втора връзка, TLS и
GSSAPI или SSPI; за MySQL двоичният протокол за подготвени заявления, компресията,
\code{caching\_sha2\_password} и отново TLS.
Декодирането обхваща едномерни масиви, а не вложени, и прехвърля отговорност за живота на
данните върху приложението. Без TLS библиотеката е годна в доверена мрежа и не е годна през
публична. Сравнение с \code{libpq} и \code{libmysqlclient} не е направено, защото те не са
налични на машината, и твърдение за отношението към тях не се прави.

\textbf{Инженерна трактовка.} Три пъти измерването поправи очакване. Разширеният цикъл
\emph{не} струва допълнително обръщане при първо изпълнение, защото петте съобщения се
изпращат в един запис. Измерването откри и дефект, който тестовете не бяха уловили: при пул с
размер едно единият работник гладуваше, защото освободената връзка се връщаше в общия списък;
поправката е пряко предаване, покрита с тест. Третото е методическо: случаят с пропуск в кеша
първоначално измерваше и изместване, тоест даваше 15 съобщения и 2 обръщания, вярно за
изпълняваното и грешно спрямо етикета си.

\textbf{По-нататъшна работа.} Сравнение с доставчиковите библиотеки на подходяща машина; TLS;
трети протокол, който би проверил доколко общият слой е общ; съвместимост с CLI~\cite{sql_cli}.
